% EN48 c10 | 原书 p.131（PDF 140） | 第 6 章 记号
% 骨架文件：标题行已设定，正文由译者填写。每页起首用 % ===== 原书 p.N ===== 注释标记。
\section{记号}

% ===== 原书 p.131 =====
\medskip\noindent\textbf{1.2 节}\par
\begin{itemize}[leftmargin=2.2em,itemsep=0.25em]
\item $H=\Sigma_2$：亏格 2 曲线；
\item $V=(\C^*)^2/\Gamma_B$：阿贝尔簇，$H$ 是它的超曲面；
\item $\mathcal{L}\to V$：$V$ 上次数为 $(1,1)$ 的线丛；
\item $s:V\to\mathcal{L}$：多 theta 函数（multi-theta function），即 $\mathcal{L}$ 的截面；
\item $\lambda$：线性形式，决定 $\mathcal{L}$ 的 Chern 类；
\item $\kappa$：二次形式，决定 $\mathcal{L}$ 的全纯结构。
\end{itemize}

\medskip\noindent\textbf{1.3 节}\par
\begin{itemize}[leftmargin=2.2em,itemsep=0.25em]
\item $M$：特征格，或环簇上的函数格；
\item $N$：余特征格，或环簇的单参数子群格；
\item $M_\R=M\otimes_\Z\R$，$N$ 同理；
\item $m$：$M$ 或 $M_\R$ 的元素；
\item $u$：$N$ 中的格元素。
\end{itemize}

\medskip\noindent\textbf{1.4 节}\par
\begin{itemize}[leftmargin=2.2em,itemsep=0.25em]
\item $\widetilde{\Delta}_Y$：定义 $\widetilde{Y}$ 的多面体；
\item $\Delta_Y=\widetilde{\Delta}_Y/\Gamma_B$；
\item $(\xi_1,\xi_2,\eta)\in M_\R/\Gamma_B\cong\R^3/\Gamma_B$：表示多面体上的坐标；
\item $v_0:Y\to\C$：超势；
\item $V^\vee\cong(\C^*)^2/\Gamma_B$：$v_0$ 的纤维。
\end{itemize}

\transnote{$\Gamma_B$ 是第 1 章引入的格，作用在 $(\C^*)^2$ 上；本页是记号表，只列出记号而不重复定义。阅读此处时应回头参看第 1 章的相应约定。}

\transnote{本节的小标题「1.2 节」「1.3 节」「1.4 节」指向原书正文的章节，不是本译注的章节编号。列出的是各节所用记号的汇总，条目顺序照原文，未重新分类。}

\transnote{$\Delta_Y$ 与 $\widetilde{\Delta}_Y$ 只差一个商：$\widetilde{\Delta}_Y$ 是万有覆盖上的多面体，$\Delta_Y$ 则是取 $\Gamma_B$ 商之后的多面体。两者在原文中都用 $\Delta$ 系记号，容易混淆；凡出现商的情形，本文保留下标与 tilde 的区别。}
